\section{Two local Coulomb branches}\label{sec:local}

Put \(a_1=e_1=(1,0,0)\) and \(a_2=e_2=(0,1,0)\).
At each of the two triples \((0,a_k,-a_k)\),
\begin{equation}\label{eq:center-jets}
 c(0,a_k,-a_k)=\frac52,\qquad
 \nabla_xc(0,a_k,-a_k)=0,\qquad
 \nabla_zc(0,a_k,-a_k)=\frac54a_k.
\end{equation}
The equal cost values and vanishing central gradients allow the two
families to share one central potential.

Near each center triple, we use the positive outer point \(y\) as the parameter and find
the central and negative outer points by minimization. Subtracting
suitable concave quadratic potentials in those two coordinates makes
the minimum unique; its value defines the potential at \(y\).
We also require both minimizing coordinates to depend diffeomorphically
on \(y\). We can then invert the central coordinate and obtain both
outer components as diffeomorphic images of a common central set.

\begin{lemma}\label{lem:local}
There are a number \(K>0\), open neighborhoods \(P_0^k,P_k^+,P_k^-\)
of \(0,a_k,-a_k\), respectively, and smooth maps
\(X_k,Z_k:P_k^+\to\R^3\), for \(k=1,2\), with the following
properties:
\begin{enumerate}
 \item Each of \(X_k\) and \(Z_k\) is a diffeomorphism onto its
 open image, and
 \[
  X_k(a_k)=0,\qquad Z_k(a_k)=-a_k.
 \]
 \item Define
 \[
 u_0(x)=\frac52-\frac K2|x|^2,\qquad
 v_k(z)=\frac54a_k\cdot(z+a_k)-\frac K2|z+a_k|^2.
 \]
 There is a smooth function \(w_k:P_k^+\to\R\), with
 \(w_k(a_k)=0\), such that on
 \(P_0^k\times P_k^+\times P_k^-\),
 \begin{equation}\label{eq:local-certificate}
 c(x,y,z)\ge u_0(x)+w_k(y)+v_k(z).
 \end{equation}
 Equality holds if and only if
 \((x,z)=(X_k(y),Z_k(y))\).
\end{enumerate}
\end{lemma}

\begin{proof}
For a parameter \(K>0\), use the potentials in the statement and set
\[
 F_k(x,y,z)=c(x,y,z)-u_0(x)-v_k(z).
\]
By \eqref{eq:center-jets}, its \((x,z)\)-gradient vanishes at the
\(k\)-th center triple. We choose \(K\) to obtain a unique minimizing
pair on a small product neighborhood for each fixed \(y\), and to make each coordinate of the
resulting minimizing pair an invertible function of \(y\).

The Coulomb cost is smooth near both center triples. At the
\(k\)-th center, write
\[
 M_k=D^2_{(x,z)(x,z)}c,\qquad
 N_k=D_y\nabla_{(x,z)}c.
\]
The \((x,z)\)-Hessian of \(F_k\) there is \(M_k+KI_6\), while its
mixed derivative in \(y\) is \(N_k\). For \(h(b)=|b|^{-1}\), \(b\ne0\),
differentiation gives
\begin{equation}\label{eq:pair-hessian}
 D^2h(b)=|b|^{-3}
 \left(3\frac{bb^{\mathsf T}}{|b|^2}-I_3\right).
\end{equation}
This matrix has eigenvalues \(2|b|^{-3},-|b|^{-3},-|b|^{-3}\)
and is therefore invertible. The two \(3\times3\) blocks of \(N_k\) are
\[
 -D^2h(-a_k),\qquad -D^2h(-2a_k),
\]
and hence are both invertible. The matrices \(M_k\) are fixed
symmetric \(6\times6\) matrices. For all sufficiently large \(K\),
both \(M_k+KI_6\) are positive definite. Moreover,
\[
 K(M_k+KI_6)^{-1}N_k
   =(I_6+K^{-1}M_k)^{-1}N_k\longrightarrow N_k.
\]
Invertibility is an open condition on \(3\times3\) matrices, so
we may choose a single \(K\) for which both blocks of
\((M_k+KI_6)^{-1}N_k\) are invertible for each \(k\).

Fix such a \(K\). The implicit function theorem applied to
\(\nabla_{(x,z)}F_k=0\) gives a smooth stationary branch
\[
 \nabla_{(x,z)}F_k(X_k(y),y,Z_k(y))=0,
\]
with the prescribed center values. Differentiating this equation
at \(y=a_k\) gives
\[
 D(X_k,Z_k)(a_k)=-(M_k+KI_6)^{-1}N_k.
\]
The two block invertibilities allow the inverse function theorem
to be applied separately to \(X_k\) and \(Z_k\).
After restricting to one sufficiently small neighborhood of \(a_k\),
both maps are diffeomorphisms onto their open images.

By continuity, the \((x,z)\)-Hessian of \(F_k\) remains positive
definite on a product of sufficiently small open balls around the
three centers. Shrink the \(y\)-ball so that the stationary pair
\((X_k(y),Z_k(y))\) belongs to the product of the other two balls.
For fixed \(y\), this latter product is convex. The positive
definite Hessian makes \(F_k(\cdot,y,\cdot)\) strictly convex on it,
so its stationary pair is its unique minimum there.

Define its minimum value by
\[
 w_k(y)=F_k(X_k(y),y,Z_k(y)).
\]
The preceding minimum property gives
\eqref{eq:local-certificate} with exactly the stated equality set.
The center values in \eqref{eq:center-jets} give \(w_k(a_k)=0\).
All neighborhoods can be restricted further as needed to retain
the separate diffeomorphism properties.
\end{proof}
